\scanpage{104}
On peut être plus explicite encore

\deleted{Plaçons} Les $\delta_B$ pour $B$ variable forment une section du fibré $\underline{O}_P(\unsure{2|1})\inserted{=\underline{O}_P(1)^{\otimes2}}$ sur le fibré projectif (de dim relative $2$)
\[
\mathbb{P}(A/k)
\]
des (hyper)plans dans $A$, donc proviennent d'un élémt canonique de
\[
\begin{aligned}
\Gamma(\underline{O}_P(2))&\simeq\operatorname{Sym}^2_k(A/k)=\\
&\simeq\text{espace des formes quadratiques}\\
&\hspace{1.4em}\text{sur}\quad(A/k)^\vee\simeq A^{\unsure{s|\leftrightarrow|\natural}}
\end{aligned}
\]

Or il y a une telle forme quadratique canonique sur $A^s$, savoir $\det$. Et on vérifie que $\delta_B$ est bien défini par cet élémt.

(Cela n'a rien à voir avec la structure multiplicative de $A$, qui n'intervient pas pour définir les $\delta_B$, il suffit de la forme quadratique \inserted{non dég.} $\det$ et l'élémt $1$ de $A$ \deleted{\unsure{non isotro|non injection}} \ldots
% archive sheet number: 103
